\section{Summary}

In this paper, we have presented the results that form the basis for
the ongoing efforts to calculate gluon PDF using the pseudo-PDF approach.

In particular, we gave a classification of possible two-gluon correlator functions.
We have identified those of them that contain the invariant amplitude ${\cal M}_{pp}(\nu,-z^2)$
that determines the gluon PDF in the light-cone $z^2 \to 0$ limit.
Since this limit is singular, one needs the matching conditions
that relate ${\cal M}_{pp}(\nu,z_3^2)$ to the light-cone PDF $f(x,\mu^2)$.

To this end, using the method of Ref. \cite{Balitsky:1987bk}, we
have performed calculations of the one-loop corrections
to the gauge-invariant correlator of two gluon field-strength tensors,
with all Lorentz indices explicit.
To preserve gauge invariance, we have used the dimensional regularization.

Since the DR produces the same form $\ln z_3^2 \mu^2$ both for
logarithms related to the UV singularities
and for those reflecting the DGLAP evolution, we have made an effort to
separate these two sources of the \mbox{$\ln z_3^2$-dependence}
at small $z_3^2$.
When we form a reduced ITD ${\mathfrak M}(\nu,z_3^2)$, the UV-related contributions are canceled, and
only the DGLAP-related terms remain in the matching relation
between the reduced ITD and the light-cone ITD.

The matching relation may be also written in a kernel form (\ref{MtchI})
that directly connects lattice data on ${\mathfrak M}(\nu,z_3^2)$
with the normalized
gluon PDF $x f_g (x,\mu^2)/ {\langle x \rangle_{\mu^2}}$.
The average gluon momentum fraction ${\langle x \rangle_{\mu^2}}$
needs to be extracted from a separate lattice calculation.

We have also demonstrated that our results
may be used for a rather straightforward calculation
of the
one-loop corrections to quasi-PDFs,
providing new insights concerning their structure
that may be used to check the results
for the gluon quasi-PDF matching conditions.

In a future publication, we plan to present more details
of our calculations, and to give a complete result for the box diagram,
in particular for the non-forward kinematics
that are needed in lattice calculations
of distribution amplitudes and GPDs.
We also plan to include calculations for gluon-quark
and quark-gluon terms.
